\documentclass[11pt]{article}
\usepackage{amsmath,amssymb,amsthm}
\usepackage[margin=1in]{geometry}
\usepackage{hyperref}

\newtheorem{theorem}{Theorem}
\newtheorem{lemma}{Lemma}
\theoremstyle{definition}
\newtheorem{definition}{Definition}

\title{Exact $\mathfrak{sl}(n,\mathbb{R})$ from Trigonometric Operators:\\
$\mathfrak{sl}(3)$, $\mathfrak{sl}(4)$, $\mathfrak{sl}(5)$, $\mathfrak{sl}(6)$}
\author{Kit Cosby}
\date{October 11, 2026}

\begin{document}
\maketitle

\begin{abstract}
From six trigonometric operators ($srx$, $sxp$, $cxp$, $crx$, $urx$, $uxp$) we construct exact matrix families that generate $\mathfrak{sl}(3,\mathbb{R})$, $\mathfrak{sl}(4,\mathbb{R})$, $\mathfrak{sl}(5,\mathbb{R})$, and $\mathfrak{sl}(6,\mathbb{R})$ over exact algebraic arithmetic ($\overline{\mathbb{Q}}$). Each identification is proved via exact rank and exact Lie closure; for $n=3,4,5$ the Killing signature provides independent confirmation, and for $n=6$ the identification follows by dimension count. All results are scope-labeled: PROVED (exact), CHECKED (numeric), ASSERTED (unproved).
\end{abstract}

\section{Operators and the base matrix}

\begin{definition}[Core operators]
For $x$ in Quadrant I:
\begin{align*}
srx &= |\csc x| + \cot x = \frac{1+\cos x}{\sin x}, &
sxp &= 1/srx, \\
cxp &= |\sec x| + \tan x = \frac{1+\sin x}{\cos x}, &
crx &= 1/cxp, \\
urx &= srx - crx, &
uxp &= cxp - sxp.
\end{align*}
\end{definition}

\begin{definition}[Base matrix]
\[
T_a(x) =
\begin{pmatrix}
urx & srx & 0 \\
cxp & -urx & uxp \\
0 & -uxp & 0
\end{pmatrix}.
\]
\end{definition}

\begin{theorem}[Double-angle identities, PROVED]
\[
\frac{\sin(2x)}{4} = \frac{1}{urx+uxp}, \qquad
\cos(2x) = \frac{srx+crx-sxp-cxp}{urx+uxp}.
\]
\end{theorem}

\begin{proof}
Set $t=\tan(x/2)$. Then $\sin x = \frac{2t}{1+t^2}$, $\cos x = \frac{1-t^2}{1+t^2}$, and:
\[
srx = \frac{1}{t}, \quad sxp = t, \quad
cxp = \frac{1+t}{1-t}, \quad crx = \frac{1-t}{1+t}.
\]
Compute:
\[
urx = srx-crx = \frac{1}{t}-\frac{1-t}{1+t} = \frac{1+t^2}{t(1+t)},
\qquad
uxp = cxp-sxp = \frac{1+t}{1-t}-t = \frac{1+t^2}{1-t}.
\]
Then:
\[
urx+uxp = (1+t^2)\left(\frac{1}{t(1+t)}+\frac{1}{1-t}\right)
= \frac{(1+t^2)^2}{t(1-t^2)}.
\]
Now $\sin(2x)/4 = \frac{1}{2}\sin x\cos x = \frac{t(1-t^2)}{(1+t^2)^2} = \frac{1}{urx+uxp}$, proving the first identity.

For the second: $srx+crx = \frac{1+2t-t^2}{t(1+t)}$, $sxp+cxp = \frac{1+2t-t^2}{1-t}$, so:
\[
srx+crx-sxp-cxp = \frac{(1+2t-t^2)(1-2t-t^2)}{t(1-t^2)}.
\]
Dividing by $urx+uxp = \frac{(1+t^2)^2}{t(1-t^2)}$ gives $\frac{(1+2t-t^2)(1-2t-t^2)}{(1+t^2)^2} = \frac{1-6t^2+t^4}{(1+t^2)^2} = \cos(2x)$, since $\cos(2x)=\cos^2x-\sin^2x=\frac{(1-t^2)^2-4t^2}{(1+t^2)^2}$.
\end{proof}

\section{Exact $\mathfrak{sl}(3,\mathbb{R})$}

\begin{theorem}[PROVED]
Eight exact traceless $3\times3$ matrices derived from $T_a(x)$ at $x\in\{15^\circ,30^\circ,45^\circ,60^\circ\}$ plus commutators:
\begin{itemize}
\item have rank 8 (exact, over $\overline{\mathbb{Q}}$);
\item close under all $\binom{8}{2}=28$ commutators (exact);
\item have Killing signature $(5,3)$ (exact);
\item have quadratic Casimir $C_2 = \frac{4}{9}I_3$ (exact);
\item have 2-dimensional centralizer (Cartan);
\item have $A_2$ root system (six nonzero roots in $\pm$ pairs).
\end{itemize}
Therefore they span $\mathfrak{sl}(3,\mathbb{R})$.
\end{theorem}

\noindent\textit{Proof scripts:} E70 (rank+closure), E71 (Casimir), E72 (roots/Cartan). All computations exact over $\overline{\mathbb{Q}}$. \qed

\section{Exact $\mathfrak{sl}(4,\mathbb{R})$}

\begin{definition}[Block extension]
\[
X_4(x) =
\begin{pmatrix}
T_a(x) & v(x) \\
v(x)^T & 0
\end{pmatrix},
\qquad
v(x) =
\begin{pmatrix}
\sin(2x)/4 \\
\cos(2x) \\
0
\end{pmatrix},
\]
where the components are given by the Lemma via $urx,uxp,srx,crx,sxp,cxp$.
\end{definition}

\begin{theorem}[PROVED]
Fifteen exact traceless $4\times4$ matrices from $X_4(x)$:
\begin{itemize}
\item have rank 15 (exact);
\item close under all $\binom{15}{2}=105$ commutators (exact);
\item have Killing signature $(9,6,0)$ (exact).
\end{itemize}
Since 15 independent traceless $4\times4$ matrices span the full $\mathfrak{sl}(4,\mathbb{R})$, and the Killing signature confirms the split real form, this is $\mathfrak{sl}(4,\mathbb{R})$.
\end{theorem}

\noindent\textit{Proof script:} E88. \qed

\section{Exact $\mathfrak{sl}(5,\mathbb{R})$}

\begin{definition}[Block extension]
\[
X_5(x) =
\begin{pmatrix}
X_4(x) & v_4(x) \\
v_4(x)^T & 0
\end{pmatrix},
\qquad
v_4(x) =
\begin{pmatrix}
\sin(4x) \\
\cos(4x) \\
0 \\
0
\end{pmatrix},
\]
where $\sin(4x)=8S_2C_2$ and $\cos(4x)=C_2^2-16S_2^2$ with $S_2=\sin(2x)/4$, $C_2=\cos(2x)$.
\end{definition}

\begin{theorem}[PROVED]
Twenty-four exact traceless $5\times5$ matrices from $X_5(x)$:
\begin{itemize}
\item have rank 24 (exact);
\item close under all $\binom{24}{2}=276$ commutators (exact);
\item have Killing signature $(14,10,0)$ (exact).
\end{itemize}
Therefore they span $\mathfrak{sl}(5,\mathbb{R})$.
\end{theorem}

\noindent\textit{Proof scripts:} E90b (rank+closure), E90c (Killing). \qed

\section{The $\mathfrak{sl}(6)$ boundary}

\begin{definition}[Block extension]
\[
X_6(x) =
\begin{pmatrix}
X_5(x) & v_5(x) \\
v_5(x)^T & 0
\end{pmatrix},
\qquad
v_5(x) =
\begin{pmatrix}
\sin(8x) \\
\cos(8x) \\
0 \\
0 \\
0
\end{pmatrix}.
\]
\end{definition}

\begin{theorem}[PROVED]
Thirty-five exact traceless $6\times6$ matrices from $X_6(x)$:
\begin{itemize}
\item have rank 35 (exact, confirmed by independent re-verification);
\item close under all $\binom{35}{2}=595$ commutators (exact).
\end{itemize}
Therefore they form a 35-dimensional Lie subalgebra of $\mathfrak{sl}(6,\mathbb{R})$. Since $\dim\mathfrak{sl}(6,\mathbb{R})=35$, this subalgebra is the whole algebra. Therefore it is $\mathfrak{sl}(6,\mathbb{R})$.
\end{theorem}

\noindent\textit{Proof scripts:} E91 (rank+closure), E94 (independent rank confirmation). The identification follows by dimension count; no Killing signature is needed. An earlier numeric Killing computation giving $(18,17,0)$ is discarded as numerically unstable. \qed

\begin{corollary}
The exact block-extension ladder yields $\mathfrak{sl}(3,\mathbb{R})$, $\mathfrak{sl}(4,\mathbb{R})$, $\mathfrak{sl}(5,\mathbb{R})$, and $\mathfrak{sl}(6,\mathbb{R})$.
\end{corollary}

\section{Scope and limitations}

\begin{itemize}
\item \textbf{PROVED:} Double-angle identities (analytic). $\mathfrak{sl}(3,\mathbb{R})$, $\mathfrak{sl}(4,\mathbb{R})$, $\mathfrak{sl}(5,\mathbb{R})$ via exact rank, closure, Killing signature. $\mathfrak{sl}(6,\mathbb{R})$ via exact rank and closure plus dimension count.
\item \textbf{CHECKED:} None (the earlier numeric $(18,17,0)$ Killing for $n=6$ is discarded as numerically unstable).
\item \textbf{ASSERTED:} None in this paper.
\item \textbf{OPEN:} Whether the ladder extends to $\mathfrak{sl}(7,\mathbb{R})$ and beyond. A different choice of angles may avoid the $8x$ aliasing.
\end{itemize}

No physical labels (space, time, force, particle) are assigned. No connection to $E_8$ or M-theory is claimed. The $A$-series ladder is classical Lie theory; it does not derive exceptional algebras.

\section{Proof scripts}

All scripts at \url{https://github.com/kitcosby/geometry-first/tree/main/proofs}:
\begin{itemize}
\item E70, E71, E72: $\mathfrak{sl}(3,\mathbb{R})$.
\item E88: $\mathfrak{sl}(4,\mathbb{R})$.
\item E90b, E90c: $\mathfrak{sl}(5,\mathbb{R})$.
\item E91: $\mathfrak{sl}(6,\mathbb{R})$ rank+closure.
\item E94: Independent exact rank confirmation for $\mathfrak{sl}(6,\mathbb{R})$.
\item E93: $\mathfrak{sl}(4,\mathbb{R})$ with non-zero $v[2]$ (robustness).
\end{itemize}

\end{document}
